% Generated draft: requires mathematical and source review.
\documentclass[11pt]{article}
\usepackage[margin=25mm]{geometry}
\usepackage{fontspec}
\setmainfont{Latin Modern Roman}
\usepackage{amsmath,amssymb,mathtools,unicode-math}
\setmathfont{Latin Modern Math}
\usepackage{enumitem,needspace,xurl,hyperref}
\setlength{\parindent}{0pt}
\setlength{\parskip}{5pt}
\setlength{\emergencystretch}{3em}
\newcommand{\N}{\mathbb{N}}
\newcommand{\Z}{\mathbb{Z}}
\newcommand{\Q}{\mathbb{Q}}
\newcommand{\R}{\mathbb{R}}
\newcommand{\C}{\mathbb{C}}
\newcommand{\F}{\mathbb{F}}
\newcommand{\A}{\mathbb{A}}
\newcommand{\PP}{\mathbb{P}}
\newcommand{\E}{\mathbb{E}}
\newcommand{\eps}{\varepsilon}
\newcommand{\dd}{\,\mathrm{d}}
\begin{document}
\section*{UnsolvedMath ID: 4800016}
\textbf{Source problem code:} catalogue entry
\subsection{Definitions and mathematical statement}
The excerpt refers to averages labeled ``E:Multies'' without reproducing their definition. In the standard commuting-transformation interpretation, take a probability space \((X,\mathcal B,\mu)\), commuting invertible measure-preserving transformations \(T_1,\ldots,T_\ell\), and \(f_1,\ldots,f_\ell\in L^\infty(\mu)\). For integer polynomials \(p_i\), the averages are
\[
 A_N(f_1,\ldots,f_\ell)=\frac1N\sum_{n=1}^N\prod_{i=1}^\ell f_i\circ T_i^{p_i(n)}.
\]
For each transformation define its rational Kronecker factor by the completed invariant sigma-algebra
\[
 \mathcal K_{\mathrm{rat}}(T_i)=\bigvee_{q\ge1}\mathcal I(T_i^q),
 \qquad
 \mathcal I(T_i^q)=\{B\in\mathcal B:\mu(T_i^{-q}B\mathbin{\triangle}B)=0\}.
\]
Equivalently this factor is generated by eigenfunctions whose eigenvalues are roots of unity, with invariant functions included.

Under the usual polynomial convention, rational independence means independence modulo constant polynomials:
\[
 \sum_{i=1}^{\ell}c_i p_i(t)\text{ constant},\quad(c_1,\ldots,c_\ell)\in\mathbb Q^\ell
 \quad\Longrightarrow\quad c_1=\cdots=c_\ell=0.
\]
The requested characteristic-factor assertion is that for every bounded tuple,
\[
 \left\|A_N(f_1,\ldots,f_\ell)
 -A_N\bigl(\mathbb E(f_1\mid\mathcal K_{\mathrm{rat}}(T_1)),\ldots,
 \mathbb E(f_\ell\mid\mathcal K_{\mathrm{rat}}(T_\ell))\bigr)\right\|_{L^2(\mu)}
 \longrightarrow0.
\]
No ergodicity assumption is inserted. Invertibility accommodates negative polynomial iterates. This precise version is conditional on the standard interpretation: the source excerpt omits the system hypotheses, definition of rational independence, convergence mode, and the referenced equation, so they require confirmation from the original problem.

The registry labels the problem solved and attributes the assertion to F.--Kuca, \emph{Inventiones} (2025), without a full citation or exact theorem hypotheses. Its background separately says the author's November 2025 progress page does not record a resolution of this numbered problem. These supplied status reports are retained without reconciling them by an unsupported literature claim.
\subsection{Short English statement}
For commuting transformations and rationally independent polynomial iterates, replacing each function by its projection onto that transformation's rational Kronecker factor should leave the multiple averages asymptotically unchanged in \(L^2\). The omitted referenced average and independence convention require source confirmation; a supplied 2025 solved-status attribution is unverified here.
\subsection{Sources}
\begin{itemize}
\item Local registry: \url{data/unsolvedmath/problems.json}, permanent UnsolvedMath ID \(4800016\), source code catalogue entry; supplied statement, background, and research summary, updated 2026-08-06.
\item Nikos Frantzikinakis, \emph{Some open problems on multiple ergodic averages} (2016), Problem 16: \url{https://arxiv.org/abs/1103.3808}. Supplied source-TeX extraction accessed 2026-07-29. The background reports no resolution entry for this problem on the author's November 2025 progress page; no separate page URL is supplied.
\item F.--Kuca, \emph{Inventiones} (2025), as abbreviated in the supplied summary; no title, volume, pages, URL, or exact theorem statement is provided.
\end{itemize}
\end{document}
